\documentclass[11pt,a4paper]{article}
\usepackage[utf8]{inputenc}
\usepackage[T1]{fontenc}
\usepackage[brazilian]{babel}
\usepackage[margin=2.3cm]{geometry}
\usepackage{booktabs,array,tabularx}
\usepackage{xcolor}
\usepackage{pgfplots}
\pgfplotsset{compat=1.17}
\PassOptionsToPackage{hyphens}{url}
\usepackage[hidelinks]{hyperref}
\usepackage{enumitem}
\setlist{nosep,leftmargin=1.4em}
\emergencystretch=2em

\definecolor{impl}{HTML}{3B6EA8}
\definecolor{rev}{HTML}{E0A030}

\title{Sweatshop no PhysiSyst: corda solta da polia com massa, desenho e clique no caminho simulado\\\large Sonnet 5.5 (high, hard xhigh) no stage 2, Opus 5.5 (xhigh) no stage 3}
\author{Foreman, sessão \texttt{sweatshop/2026-10-02-1156}}
\date{2 de outubro de 2026}

\begin{document}
\maketitle

\section*{Resumo}

Run diurno em dois lançamentos na mesma sessão. O primeiro (11:56) usou o
lineup padrão de então, \texttt{Models:} com \texttt{claude-opus-5-5 high} no
stage 2 e \texttt{claude-fable-5-1 high} no stage 3. Às 12:08 o Bruno pediu
para parar depois do estágio corrente e trocar o lineup. O driver parou pelo
arquivo \texttt{STOP} às 12:13 e foi relançado às 12:15 com
\texttt{-Models 'stage 2 sonnet-5.5 high, hard sonnet-5.5 xhigh, stage 3
opus-5.5 xhigh'}, só para este run. No mesmo intervalo o \texttt{Models:} do
\texttt{AGENTS.md} virou \texttt{gpt-6.1-sol high}, hard \texttt{sonnet-5.5
xhigh}, \texttt{gpt-6.1-sol max} (\texttt{1948857} na \texttt{main},
\texttt{f42a2c9} na sessão); o antigo ficou como \texttt{Models (opus)}.

\begin{itemize}
  \item \textbf{Produzido:} 6 tickets mergeados, \textbf{nenhuma reabertura}:
        CLEAN-26, CLEAN-25, PHY-56, CLEAN-27, PHY-55 (\texttt{Difficulty:
        hard}) e PHY-57. A fila esvaziou. PR de sessão:
        \href{https://github.com/Laginho/PhysiSyst/pull/13}{\#13}, 6 linhas,
        todas \texttt{Review: agent} --- Approve.
  \item \textbf{Custo:} \textbf{\$31,40} a preço de lista nos 12 estágios,
        todos com preço: \texttt{claude-opus-5-5 high} \$2,31 (2),
        \texttt{claude-fable-5-1 high} \$1,73 (1), \texttt{sonnet-5.5 high}
        \$6,46 (3), \texttt{sonnet-5.5 xhigh} \$3,63 (1), \texttt{opus-5.5
        xhigh} \$17,27 (5). 84 min de estágio, 48,7\,M tokens de entrada
        (95,8\,\% em cache) e 372\,k de saída. \$5,23 por ticket mergeado.
  \item \textbf{O que deu errado:} nada quebrou. O ponto fraco é de medida: a
        revisão do Opus 5.5 xhigh fechou sozinha, em 4 de 5 tickets, lacunas que
        outro revisor teria devolvido (um teste faltando no CLEAN-27, o
        vermelho isolado do critério 4.3 no PHY-56). O scoreboard registra 0
        reaberturas e não vê esses achados.
  \item \textbf{Veredito provisório:} lineup caro e sem atrito. A revisão
        custou \$3,45 em média, 5 vezes a do Sol max (\$0,67 no run anterior),
        e levou 61\,\% do custo do run. O Sonnet 5.5 entregou 4 de 4 em
        \texttt{to-review} de primeira. Com 3 e 1 revisões nos pares novos, é
        anedota.
\end{itemize}

\section{Linha do tempo}

\begin{tabularx}{\linewidth}{@{}lX@{}}
\toprule
11:56 & Lançamento 1, driver \texttt{74dc512}, lineup padrão. Sessão nova
        \texttt{sweatshop/2026-10-02-1156} (o PR \#12 já estava mergeado).\\
12:01 & CLEAN-26 (um \texttt{wrapAngle} exportado em \texttt{ropePath}) em
        \texttt{to-review} em 4 min. 12:04: mergeado pelo Fable, sem achados.\\
12:08 & O Bruno pede parada depois do estágio corrente e novo lineup;
        \texttt{STOP} criado com o CLEAN-25 em implementação.\\
12:13 & CLEAN-25 (junções das polias soltas igualmente espaçadas) em
        \texttt{to-review} em 9 min. O driver acha o \texttt{STOP}, publica o
        PR \#13 com 1 ticket e volta à \texttt{main}.\\
12:15 & Lineup padrão trocado no \texttt{AGENTS.md}; lançamento 2 com
        \texttt{-Models}. A \texttt{main} tinha 4 commits do stage 1 sem push,
        que subiram junto.\\
12:22 & CLEAN-25 mergeado pelo Opus 5.5 xhigh em 7 min.\\
12:33 & PHY-56 (desenho da corda pelo caminho simulado) em 11 min. 12:43:
        mergeado; a revisão isolou o vermelho do critério 4.3.\\
12:47 & CLEAN-27 (clique na corda pelo caminho simulado) em 3 min. 12:53:
        mergeado; a revisão acrescentou o clique antes do primeiro passo.\\
13:03 & PHY-55 (corda se solta da polia com massa) vai para o hard
        (\texttt{sonnet-5.5 xhigh}): 78 testes vermelhos, \texttt{to-review}
        em 10 min. 13:11: mergeado.\\
13:14 & PHY-57 (corda volta a prender no disco) em 3 min. 13:23: mergeado.\\
13:24 & ``Nothing runnable''; PR \#13 atualizado com 6 tickets.\\
\bottomrule
\end{tabularx}

\section{Métricas por estágio}

Linhas 129--140 de \texttt{.scratch/run-log.md}. Minutos, tokens e custo como o
driver gravou.

\medskip
\noindent{\footnotesize
\begin{tabularx}{\linewidth}{@{}l l r r r r r X@{}}
\toprule
Estágio & Modelo & Min & Entrada & Cache & Saída & Custo & Desfecho\\
\midrule
C26 impl.    & opus-5-5 high  & 4  & 1,00\,M & 93\,\% & 5\,k  & 0,907 & to-review\\
C26 revisão  & fable-5-1 high & 3  & 0,91\,M & 93\,\% & 4\,k  & 1,734 & mergeado\\
C25 impl.    & opus-5-5 high  & 9  & 2,60\,M & 97\,\% & 16\,k & 1,402 & to-review\\
C25 revisão  & opus xhigh     & 7  & 5,50\,M & 96\,\% & 37\,k & 3,186 & mergeado\\
P56 impl.    & sonnet high    & 11 & 6,60\,M & 96\,\% & 47\,k & 4,158 & to-review\\
P56 revisão  & opus xhigh     & 10 & 5,10\,M & 95\,\% & 48\,k & 3,488 & mergeado\\
C27 impl.    & sonnet high    & 3  & 2,40\,M & 97\,\% & 11\,k & 1,164 & to-review\\
C27 revisão  & opus xhigh     & 6  & 4,20\,M & 96\,\% & 35\,k & 2,631 & mergeado\\
P55 impl.    & sonnet xhigh   & 10 & 5,90\,M & 97\,\% & 52\,k & 3,628 & to-review\\
P55 revisão  & opus xhigh     & 9  & 7,20\,M & 96\,\% & 53\,k & 4,069 & mergeado\\
P57 impl.    & sonnet high    & 3  & 2,10\,M & 97\,\% & 12\,k & 1,142 & to-review\\
P57 revisão  & opus xhigh     & 9  & 5,20\,M & 94\,\% & 52\,k & 3,893 & mergeado\\
\midrule
\multicolumn{2}{@{}l}{Total} & 84 & 48,71\,M & 95,8\,\% & 372\,k & 31,402 & 12 de 12 com preço\\
\multicolumn{2}{@{}l}{Implementação} & 40 & 20,60\,M & 96,5\,\% & 143\,k & 12,401 & 6 estágios\\
\multicolumn{2}{@{}l}{Revisão} & 44 & 28,11\,M & 95,4\,\% & 229\,k & 19,001 & 6 estágios\\
\multicolumn{2}{@{}l}{Retrabalho} & 0 & --- & & --- & 0 & nenhuma reabertura\\
\multicolumn{2}{@{}l}{Perdido} & 0 & --- & & --- & 0 & nenhuma linha \texttt{?}\\
\bottomrule
\end{tabularx}}

\bigskip
\begin{center}
\begin{tikzpicture}
\begin{axis}[
  ybar stacked, bar width=18pt, width=\linewidth, height=6.5cm,
  ylabel={Tokens de entrada (M)}, ymin=0,
  symbolic x coords={CLEAN-26,CLEAN-25,PHY-56,CLEAN-27,PHY-55,PHY-57},
  xtick=data, x tick label style={font=\small},
  legend style={at={(0.02,0.97)},anchor=north west,draw=none,font=\small},
  axis lines*=left, ymajorgrids, grid style={gray!25}]
\addplot[fill=impl,draw=none] coordinates
  {(CLEAN-26,1.00) (CLEAN-25,2.60) (PHY-56,6.60) (CLEAN-27,2.40) (PHY-55,5.90) (PHY-57,2.10)};
\addplot[fill=rev,draw=none] coordinates
  {(CLEAN-26,0.91) (CLEAN-25,5.50) (PHY-56,5.10) (CLEAN-27,4.20) (PHY-55,7.20) (PHY-57,5.20)};
\legend{implementação, revisão}
\end{axis}
\end{tikzpicture}
\end{center}

\paragraph{Onde foi o custo.} Por ticket: PHY-55 \$7,70 (19 min), PHY-56
\$7,65 (21 min), PHY-57 \$5,04, CLEAN-25 \$4,59, CLEAN-27 \$3,80, CLEAN-26
\$2,64. Pela primeira vez no PhysiSyst a revisão custou mais que a
implementação: \$19,00 contra \$12,40, 58\,\% da entrada e 61\,\% do custo. Em
4 dos 5 tickets revisados pelo Opus 5.5 xhigh a revisão custou mais que a
implementação; no PHY-57, 3,4 vezes. Até o Fable 5.1 high, no CLEAN-26, custou
\$1,73 contra \$0,91 do Opus que implementou.

\paragraph{O que se perdeu.} Nada em tokens: nenhuma reabertura, nenhuma linha
\texttt{?}. Em relógio, cerca de 2 min entre o \texttt{STOP} (12:13:35) e o
primeiro estágio do lançamento 2 (12:15:35). Relógio total: 88 min, de 11:56 a
13:24.

\section{Comparação}

Os runs anteriores do PhysiSyst com preço. Os quatro primeiros revisaram com
\texttt{gpt-6.1-sol max}; este, com Fable 5.1 high (1) e Opus 5.5 xhigh (5).

\medskip
\noindent{\footnotesize
\begin{tabularx}{\linewidth}{@{}X r r r r r@{}}
\toprule
& 30/09 (1--5) & 30/09 (6) & 01/10 & 02/10 noite & Este run\\
\midrule
Implementador & sol high & sol high & sonnet high & sol + sonnet xh & opus + sonnet\\
Estágios registrados & 28 & 10 & 21 & 13 & 12\\
Tickets mergeados & 10 & 2 & 7 & 5 & 6\\
Reaberturas & 1 & 3 & 1 & 1 & 0\\
Impl.: falhou ou perguntou & 4 de 15 & 0 de 5 & 0 de 8 & 0 de 6 & 0 de 6\\
Mediana da implementação & 7 min & 5 min & 6 min & 8,5 min & 6,5 min\\
Custo médio por impl. & \$0,52 & \$0,45 & \$1,74 & \$1,89 & \$2,07\\
Mediana da revisão & 10 min & 11 min & 12,5 min & 10,5 min & 8 min\\
Custo médio por revisão & \$0,59 & \$0,62 & \$0,65 & \$0,67 & \$3,17\\
Estágios perdidos (\texttt{?}) & 2 & 0 & 5 & 1 & 0\\
Custo registrado & \$14,20 & \$5,39 & \$19,13 & \$15,32 & \$31,40\\
Custo por ticket mergeado & \$1,42 & \$2,69 & \$2,73 & \$3,06 & \$5,23\\
\bottomrule
\end{tabularx}}

\medskip
O custo por ticket quase dobrou em relação ao run da noite, e a diferença é
quase toda da revisão: com a média de revisão do Sol max (\$0,67), este run
teria custado cerca de $16,42, ou $2,74 por ticket. O Sonnet 5.5 high saiu a
\$2,15 por implementação, perto dos \$1,74 de 01/10. O Sonnet 5.5 xhigh fez o
hard do run por \$3,63, contra \$8,82 do PHY-54 na noite, mas o PHY-55 chegou
mais mastigado: o proxy tinha fixado 9 decisões no stage 1.

\section{Atribuição das reaberturas}

Nenhuma rodada \texttt{reopened} neste run, logo nenhuma linha nova em
\path{.scratch/reopen-attribution.md}. Sem \emph{drip-feeding} e sem
\emph{churn}.

Três achados foram consertados pela própria revisão em vez de devolvidos.
Depois da recomendação 1, eles entram em \path{.scratch/reopen-attribution.md}
com rodada 0. O revisor foi \texttt{opus-5.5 xhigh}, o mesmo modelo do foreman,
então o GPT-6 Astra (medium, só leitura) rotulou os três. Os rótulos entraram
como ele os devolveu:

\medskip
\noindent{\small
\begin{tabularx}{\linewidth}{@{}l c l X@{}}
\toprule
Ticket & Rodada & Rótulo (Astra) & Achado e porquê\\
\midrule
CLEAN-27 & 0 & S2-explícito/teste & O teste de App nunca clicava antes do
  primeiro passo, então metade do critério 3.2 não tinha prova. A revisão
  acrescentou o clique (\texttt{9f95310}) e duas mutações. O critério pede o
  clique antes do passo.\\
PHY-56 & 0 & S2-explícito/teste & O critério 4.3 não tinha vermelho próprio:
  no teste compartilhado, a mutação 3 parava no 4.1. A revisão isolou e viu o
  4.3 falhar (\texttt{180309f}). O critério 7 pede que a mutação 3 exponha os
  dois.\\
PHY-57 & 0 & S2-explícito/código & O comentário de \texttt{Grip.loose} não
  dizia como o impacto do reengate é absorvido; só o ADR dizia. A revisão
  corrigiu (\texttt{42937d2}). O critério 4 pede isso no código.\\
\bottomrule
\end{tabularx}}

\medskip
Totais: 2 S2-explícito/teste, 1 S2-explícito/código, nenhum S1 ou S3. Os três
são do Sonnet 5.5 high. Sob um revisor que reabre, o par
\texttt{sonnet-5.5 high} $\to$ \texttt{opus-5.5 xhigh} teria 3 de 3 reabertos,
não 0 de 3. O foreman concorda com os três rótulos.

\section{Scoreboard}

Saída de \path{sweatshop/scripts/scoreboard.ps1} sobre os quatro repositórios
desta máquina com \path{.scratch/run-log.md} (PhysiSyst, SIGAA-ME, SynchroNice,
slip), já com as recomendações 1--3 aplicadas
(\href{https://github.com/Laginho/Agent-Skills/pull/3}{Agent-Skills\#3}): o log
local do PhysiSyst foi sincronizado com a cópia (178 linhas), as grafias
\texttt{claude-*} leem como o nome curto, e as três linhas de rodada 0 acima
contam na coluna nova.

\medskip
\noindent{\footnotesize
\begin{tabularx}{\linewidth}{@{}X r r r r r r@{}}
\toprule
Implementador & Estágios & To review & Falhou & Perguntou & Mediana até review & Custo\\
\midrule
sonnet & 121 & 75 & 31 & 15 & 7m & \$10.53 (4/121)\\
opus-5.5 & 35 & 31 & 4 & 0 & 6m & \$0.72 (1/35)\\
gpt-6.1-sol high & 35 & 31 & 0 & 4 & 9m & \$27.53 (35/35)\\
sonnet-5.5 high & 11 & 11 & 0 & 0 & 6m & \$20.42 (11/11)\\
sonnet-5.5 xhigh & 3 & 3 & 0 & 0 & 10m & \$12.44 (3/3)\\
opus-5.5 high & 2 & 2 & 0 & 0 & 9m & \$2.31 (2/2)\\
fable & 5 & 5 & 0 & 0 & 3m & ?\\
gpt-6-luna & 23 & 11 & 10 & 2 & 26m & ?\\
gpt-6-sol & 6 & 5 & 1 & 0 & 17m & ?\\
gpt-5.6-sol & 5 & 4 & 1 & 0 & 17m & ?\\
gpt-6-astra & 18 & 11 & 6 & 1 & 8m & ?\\
sonnet-5 xhigh & 11 & 6 & 2 & 3 & 8m & \$18.09 (11/11)\\
gpt-6-astra low & 2 & 2 & 0 & 0 & 30m & \$19.03 (2/2)\\
gpt-6.1-sol max & 3 & 1 & 2 & 0 & 17m & \$1.44 (1/3)\\
\bottomrule
\end{tabularx}}

\medskip
\noindent{\footnotesize
\begin{tabularx}{\linewidth}{@{}l X r r r r r r r@{}}
\toprule
Implementador & Revisor & Reviews & Reab. & Taxa & Por S2 & Taxa S2 & Sem rót. & Cons.\ S2\\
\midrule
sonnet & opus & 68 & 20 & 29\,\% & 1 & $\geq$1\,\% & 19 & 0\\
? & opus & 4 & 0 & 0\,\% & 0 & 0\,\% & 0 & 0\\
opus-5.5 & fable-5.1 & 30 & 1 & 3\,\% & 0 & $\geq$0\,\% & 1 & 0\\
gpt-6.1-sol high & gpt-6.1-sol max & 30 & 8 & 27\,\% & 6 & $\geq$20\,\% & 1 & 0\\
sonnet-5.5 high & gpt-6.1-sol max & 8 & 1 & 12\,\% & 1 & 12\,\% & 0 & 0\\
sonnet-5.5 xhigh & gpt-6.1-sol max & 2 & 1 & 50\,\% & 1 & 50\,\% & 0 & 0\\
opus-5.5 high & fable-5.1 high & 1 & 0 & 0\,\% & 0 & 0\,\% & 0 & 0\\
opus-5.5 high & opus-5.5 xhigh & 1 & 0 & 0\,\% & 0 & 0\,\% & 0 & 0\\
sonnet-5.5 high & opus-5.5 xhigh & 3 & 0 & 0\,\% & 0 & 0\,\% & 0 & 3\\
sonnet-5.5 xhigh & opus-5.5 xhigh & 1 & 0 & 0\,\% & 0 & 0\,\% & 0 & 0\\
sonnet & fable & 1 & 0 & 0\,\% & 0 & 0\,\% & 0 & 0\\
fable & fable & 4 & 1 & 25\,\% & 0 & $\geq$0\,\% & 1 & 0\\
gpt-6-luna & gpt-6-sol & 11 & 7 & 64\,\% & 0 & $\geq$0\,\% & 7 & 0\\
gpt-6-sol & gpt-5.6-sol & 5 & 4 & 80\,\% & 0 & $\geq$0\,\% & 4 & 0\\
gpt-5.6-sol & gpt-6-astra & 4 & 3 & 75\,\% & 0 & $\geq$0\,\% & 3 & 0\\
gpt-6-astra & gpt-5.6-sol & 11 & 5 & 45\,\% & 0 & $\geq$0\,\% & 5 & 0\\
sonnet-5 xhigh & opus-5.5 xhigh & 5 & 0 & 0\,\% & 0 & 0\,\% & 0 & 0\\
gpt-6-astra low & gpt-6.1-sol max & 2 & 1 & 50\,\% & 1 & 50\,\% & 0 & 0\\
gpt-6.1-sol max & gpt-6.1-sol max & 1 & 1 & 100\,\% & 0 & $\geq$0\,\% & 1 & 0\\
\bottomrule
\end{tabularx}}

\smallskip
{\footnotesize \emph{Cons.\ S2} é a coluna nova \texttt{Fixed in review for S2}:
merges cuja revisão consertou ela mesma um achado S2 (rodada 0).}

\medskip
\noindent{\footnotesize
\begin{tabularx}{\linewidth}{@{}X r r r r r r@{}}
\toprule
Implementador $\to$ Revisor & Rodadas & S2 expl.\ código & S2 expl.\ teste & S2 implícito & S1 & S3 ruído\\
\midrule
gpt-6.1-sol high $\to$ gpt-6.1-sol max & 7 & 3 & 0 & 7 & 1 & 0\\
sonnet-5.5 high $\to$ gpt-6.1-sol max & 1 & 0 & 0 & 2 & 3 & 0\\
sonnet-5.5 xhigh $\to$ gpt-6.1-sol max & 1 & 1 & 0 & 1 & 2 & 0\\
sonnet-5.5 high $\to$ opus-5.5 xhigh & 3 & 1 & 2 & 0 & 0 & 0\\
sonnet $\to$ opus & 1 & 1 & 1 & 0 & 0 & 0\\
gpt-6-astra low $\to$ gpt-6.1-sol max & 1 & 2 & 0 & 1 & 0 & 0\\
\bottomrule
\end{tabularx}}

\paragraph{O que mudou desde o relatório anterior.}
\begin{itemize}
  \item \textbf{Grafia normalizada:} \texttt{claude-opus-5-5} virou
        \texttt{opus-5.5} e \texttt{claude-fable-5-1} virou
        \texttt{fable-5.1}, nas linhas antigas e nas novas. Os números dessas
        linhas não mudaram, só o nome.
  \item \textbf{Quatro pares novos, todos deste run, com 0 reaberturas:}
        \texttt{sonnet-5.5 high} $\to$ \texttt{opus-5.5 xhigh} (3),
        \texttt{sonnet-5.5 xhigh} $\to$ \texttt{opus-5.5 xhigh} (1),
        \texttt{opus-5.5 high} $\to$ \texttt{fable-5.1 high} (1) e
        \texttt{opus-5.5 high} $\to$ \texttt{opus-5.5 xhigh} (1). O primeiro
        tem 3 consertos S2 na revisão e 3 rodadas na tabela de achados (rodada
        0).
  \item \textbf{\texttt{sonnet-5.5 high}} foi de 8 para 11 estágios e
        \textbf{\texttt{sonnet-5.5 xhigh}} de 2 para 3 (mediana de 19 para
        10 min). \texttt{opus-5.5 high} é linha nova, separada de
        \texttt{opus-5.5} sem esforço.
  \item \textbf{\texttt{gpt-6.1-sol high} $\to$ \texttt{gpt-6.1-sol max}} foi
        de 28 revisões e 21\,\% de taxa S2 para 30 e $\geq$20\,\%, com uma
        reabertura sem rótulo. \textbf{\texttt{gpt-6.1-sol max}} apareceu como
        implementador (3 estágios, 2 falhas) com um par
        \texttt{gpt-6.1-sol max} $\to$ \texttt{gpt-6.1-sol max} reaberto e sem
        rótulo. Nada disso é deste run: vem dos repositórios irmãos.
\end{itemize}

\paragraph{Como ler.}
\begin{itemize}
  \item \textbf{Mesmo implementador, dois revisores:} o Sonnet 5.5 high tem
        12\,\% de taxa S2 sob o Sol max (8 revisões) e 0\,\% sob o Opus 5.5
        xhigh (3), mas com 3 consertos S2 em 3 revisões. Somando os dois, os
        deslizes do implementador sob o Opus foram 3 de 3. \textbf{Ambos os
        pares abaixo de 10 revisões: anedota.}
  \item \textbf{Os pares do Opus 5.5 xhigh com 0\,\%} somam 10 revisões entre
        \texttt{sonnet-5 xhigh}, \texttt{sonnet-5.5 high}, \texttt{sonnet-5.5
        xhigh} e \texttt{opus-5.5 high}, nenhuma reaberta. Antes deste run a
        coluna de consertos não existia, então os 5 do \texttt{sonnet-5
        xhigh} podem esconder consertos que ninguém rotulou.
  \item \textbf{Os três rótulos novos são S2-explícito:} o ticket dizia, e o
        stage 2 não fez. Nada aponta para o stage 1.
\end{itemize}

\section{Observações sobre os modelos}

\subsection*{Opus 5.5 (high), implementação (lançamento 1)}
\begin{itemize}
  \item \textbf{2 de 2 em \texttt{to-review}, 2 de 2 mergeados.} CLEAN-26 em
        4 min (\$0,91), CLEAN-25 em 9 min (\$1,40).
  \item \textbf{Registro de mutação completo} no CLEAN-25, inclusive a mutação
        de integração nos dois cenários de aceitação.
  \item \textbf{Dois deslizes admitidos no próprio \texttt{.final}:} o
        \texttt{to-review} entrou num commit de docs separado, e o vermelho do
        critério 1 nunca mostra os 21\,m previstos porque o teste de caminhada
        para na primeira falha. Ergueu o timeout desse teste para 30\,s num
        commit só de teste e repetiu o vermelho.
\end{itemize}

\subsection*{Sonnet 5.5 (high), implementação}
\begin{itemize}
  \item \textbf{3 de 3 em \texttt{to-review}, 3 de 3 mergeados.} PHY-56 em
        11 min (\$4,16), CLEAN-27 e PHY-57 em 3 min (\$1,16 e \$1,14).
  \item \textbf{Harness errado no commit vermelho do PHY-56:} um teste nunca
        poderia passar e outro deixava uma mutação viva. Ambos foram corrigidos
        em commits só de teste, com as cinco mutações repetidas. É o mesmo
        padrão do PHY-53 do Sol high na noite.
  \item \textbf{Três lacunas de prova}, todas fechadas pela revisão (seção 4).
        Nenhuma de comportamento.
  \item \textbf{PHY-57: o primeiro commit de código deixou o
        \texttt{Stage:} em \texttt{implementing}}; um commit seguinte o
        corrigiu. Admitido no \texttt{.final}.
\end{itemize}

\subsection*{Sonnet 5.5 (xhigh), implementação hard}
\begin{itemize}
  \item \textbf{PHY-55 em 10 min, \$3,63, 78 de 78 vermelhos} no commit de
        testes (+231\,J no caso 1/1 com $v_x = 8$ antes do conserto).
  \item \textbf{Registrou e não investigou} a divergência entre as faixas de
        energia da mutação 1 e as do ticket (até +189\,kJ contra +45,9\,kJ
        previstos). A revisão explicou: o conserto também solta no
        \texttt{correctPieces}, o que o protótipo provavelmente não fazia.
  \item \textbf{Mais barato que no PHY-54} (\$3,63 contra \$8,82 com a
        reabertura), num ticket com 9 decisões do proxy já no stage 1.
\end{itemize}

\subsection*{Fable 5.1 (high), revisão (lançamento 1)}
\begin{itemize}
  \item \textbf{Uma revisão, CLEAN-26, 3 min, \$1,73:} sem achados, vermelho
        reproduzido (1 de 21). Custou 1,9 vez a implementação de um refactor de
        uma função.
\end{itemize}

\subsection*{Opus 5.5 (xhigh), revisão}
\begin{itemize}
  \item \textbf{5 de 5 mergeados, 6--10 min, \$2,63--4,07.} Repetiu o vermelho
        e as mutações em todos e rodou o gate antes de cada merge.
  \item \textbf{Conserta em vez de devolver.} CLEAN-27: escreveu o teste que
        faltava e duas mutações. PHY-56: isolou o vermelho do 4.3. PHY-55:
        explicou a divergência da mutação 1 e refluiu dois comentários que só
        passavam no \texttt{git grep} do critério 6 pela quebra de linha.
        PHY-57: cumpriu o critério 4, tirou uma linha morta e evitou recalcular
        as peças a cada passo.
  \item \textbf{Anota sem abrir ticket:} a busca ``caminho do simulador, senão
        o do documento'' em três arquivos (\texttt{draw.ts},
        \texttt{overlay.ts}, \texttt{hitTest.ts}), \texttt{drawScene} com 8
        argumentos posicionais, três regras do CLEAN-25 sem teste.
  \item \textbf{Caro:} 45\,k de saída em média por revisão, contra 13\,k do Sol
        max na noite.
\end{itemize}

\section{Driver e runtime}
\begin{itemize}
  \item \textbf{\texttt{STOP} funcionou como prometido:} achado às 12:13:35,
        entre estágios, com o PR publicado e a árvore na \texttt{main}. Foi o
        primeiro uso real; custou cerca de 2 min.
  \item \textbf{Mesmo modelo, duas grafias:} a binding escreve
        \texttt{claude-opus-5-5} e \texttt{claude-fable-5-1}; o
        \texttt{-Models} aceita \texttt{opus-5.5}. O log grava como foi
        escrito, e o scoreboard trata \texttt{claude-opus-5-5 high} e
        \texttt{opus-5.5 xhigh} como modelos diferentes. Aqui o esforço já os
        separa, mas um \texttt{opus-5.5 high} viraria uma terceira linha.
  \item \textbf{Base à frente da origem:} a \texttt{main} tinha 4 commits do
        stage 1 sem push. O Preflight aceitou (o \texttt{pull --ff-only} não
        falha com a base à frente), e a sessão nasceu deles e os publicou.
        Inofensivo aqui, mas o PR de sessão mostraria commits que a
        \texttt{main} remota ainda não tem.
  \item \textbf{Aviso inofensivo de novo:} o \texttt{.err} dos estágios do
        Sonnet traz o aviso \texttt{unrecognized\_model} do Claude Code
        para \texttt{claude-sonnet-5-5}; todos rodaram normalmente.
  \item \textbf{Log local incompleto:} \path{.scratch/run-log.md} desta máquina
        tem 140 linhas; \path{docs/run-log/run-log.md} tinha 166. Copiar o
        local por cima, como o foreman manda, apagaria 38 linhas de 30/09. O
        foreman acrescentou as 12 linhas deste run à cópia.
\end{itemize}

\section{Recomendações}
\begin{enumerate}
  \item \textbf{Registrar os consertos da revisão que acrescentam prova.}
        Quando o stage 3 escreve um teste ou fecha um critério em vez de
        reabrir, uma linha em \path{reopen-attribution.md} (com rodada 0, por
        exemplo) deixaria o scoreboard ver o achado. Hoje os pares do Opus 5.5
        xhigh mostram 0\,\% em 10 revisões somadas, e pelo menos 1 dos 3 do
        \texttt{sonnet-5.5 high} $\to$ \texttt{opus-5.5 xhigh} teria reaberto
        sob o Sol max. Linhas: \texttt{sonnet-5.5 high} $\to$
        \texttt{opus-5.5 xhigh} (3) e $\to$ \texttt{gpt-6.1-sol max} (8);
        ambas anedota.
  \item \textbf{Normalizar a grafia dos modelos no log} (\texttt{claude-opus-5-5}
        e \texttt{opus-5.5} como um nome), antes que dois registros do mesmo
        modelo e esforço se dividam no scoreboard.
  \item \textbf{Run log: mesclar, não copiar.} Ou sincronizar o
        \path{.scratch/run-log.md} desta máquina a partir de
        \path{docs/run-log/run-log.md} antes do próximo run, ou mudar o
        foreman para acrescentar as linhas do run à cópia.
  \item \textbf{O lineup padrão novo (Sol high / Sonnet 5.5 xhigh / Sol max)
        está coerente com o custo:} a revisão do Opus 5.5 xhigh custou 5 vezes
        a do Sol max sem diferença de qualidade mensurável no scoreboard.
        Apoia-se em custo, não em taxa: os pares comparáveis têm 3 e 8
        revisões.
  \item \textbf{Abrir um CLEAN para a busca de caminho repetida} em
        \texttt{draw.ts}, \texttt{overlay.ts} e \texttt{hitTest.ts}, e decidir
        se \texttt{drawScene} merece um objeto de opções. A revisão do CLEAN-27
        e a do PHY-56 anotaram; nenhum ticket existe.
  \item \textbf{Testes para as três regras do CLEAN-25 sem critério} (polia
        solta única no meio, espaçamento entre polias ainda presas, ângulo
        \texttt{start} inalterado), num CLEAN de testes ou no próximo ticket
        que mexer em \texttt{ropePath}.
  \item \textbf{A fila acabou:} nenhum ticket aberto. O próximo run depende do
        stage 1.
\end{enumerate}

Do relatório anterior: a recomendação 6 (CLEAN do \texttt{wrapAngle}) virou o
CLEAN-26 e foi mergeada aqui; a 7 (stage 1 de CLEAN-25, PHY-55 e PHY-56)
também foi feita. As outras (Primary files, proxy antes de reabrir por
fronteira, Codex por janela) não foram exercitadas neste run.

\section{Decisões do proxy}

As 31 linhas \texttt{Proxy decided} dos tickets deste run, todas de 2026-10-02 e
todas do stage 1: são as decisões tomadas em nome do Bruno. Ficam em
\texttt{\#\# Comments} de cada ticket; aqui, citadas na íntegra. CLEAN-26 e
CLEAN-27 não têm nenhuma (o CLEAN-27 nasceu de uma decisão do PHY-56).

{\small
\subsection*{CLEAN-25 (8)}
\begin{enumerate}
  \item Junções em $t = k/(n + 1)$ na ordem de \texttt{via}, sem projeção --- a
        junção não faz trabalho físico, e (d) mantém todo segmento
        $\geq |\mathrm{perna}|/(n + 1)$, o motivo original do PHY-54 para o limite.
  \item O caso do \texttt{ponytail:} (duas projeções presas na mesma ponta)
        entra neste ticket, na segunda disposição do critério 2 --- a mesma
        regra o cobre sem código a mais.
  \item Substituir a decisão do PHY-54 (ponto mais perto do centro, limitado a
        $[\frac14, \frac34]$) e corrigir o ADR --- a junção não faz trabalho
        físico, e (d) mantém todo segmento $\geq |\mathrm{perna}|/(n + 1)$, o
        motivo original do limite.
  \item \texttt{arcs[i].start} de uma polia solta fica o ângulo até a junção,
        sem mudança --- só o código de grip (sem histórico) e o
        \texttt{scenePath} (sem histórico) leem \texttt{start}.
  \item Testes só na geometria, chamando \texttt{ropePath} direto, sem cena no
        \texttt{acceptance.test.ts} --- \texttt{ropeFrame} só soma as direções
        dos segmentos, e a cena seria artificial. (Revista pela decisão 8.)
  \item O critério 5 pede o registro do mutante (b) --- só uma janela em torno
        do cruzamento separa (b) de (d).
  \item \texttt{Blocked by: CLEAN-26}, \texttt{Difficulty: normal},
        \texttt{Review: agent}, um ticket só --- mesmo arquivo do CLEAN-26,
        função privada, critérios independentes.
  \item ``Testes só na geometria'' revisto pela evidência nova (a cena simples
        da mesa com duas polias ideais já ganha energia na main): o critério de
        energia, com \texttt{twoPulleyTableScene} em
        \texttt{acceptance.test.ts} e o registro de mutação, fica neste ticket
        --- um ticket posterior nunca conseguiria ficar vermelho primeiro.
\end{enumerate}

\subsection*{PHY-56 (9)}
\begin{enumerate}
  \item A costura é \texttt{RopeState.path}, o caminho que o simulador
        resolveu na última leitura das poses reais --- desenho = simulação por
        construção, e o \texttt{readConstraints} já viaja até o \texttt{paint}
        a cada quadro; recalcular com \texttt{keep}/\texttt{sweeps} repetiria a
        regra do grip, e um \texttt{readRopePaths()} pediria outro ref no App.
  \item \texttt{path} opcional no tipo \texttt{RopeState} --- obrigatório
        quebraria o typecheck do helper \texttt{ropeState} de
        \texttt{overlay.test.ts} e de dois literais de \texttt{App.test.ts}; o
        fallback para \texttt{scenePath} existe de qualquer jeito para o
        editor.
  \item A guarda do editor fica no \texttt{paint}, numa variável só que
        alimenta \texttt{drawScene} e \texttt{tensionArrows} (\texttt{[]} sem
        \texttt{states}), e \texttt{elasticArrows} fica sem guarda --- o
        \texttt{constraintsRef} só é atualizado no boot, no reset e no rebuild;
        uma leitura velha com \texttt{via} mais curto faz
        \texttt{tensionArrows} ler um segmento indefinido e o \texttt{paint}
        lançar, e $F_{el}$ aparece em $t = 0$.
  \item \texttt{drawScene} e \texttt{tensionArrows} entram os dois --- sem as
        setas, com a corda solta, a seta em \texttt{a} aponta para o aro
        enquanto a corda desenhada vai reta.
  \item O clique na corda (\texttt{ropeAtPoint}) fica fora e vai para o
        CLEAN-27 --- é outra costura (editor, não render), e selecionar a corda
        solta durante o playback é raro.
  \item \texttt{Blocked by: none} --- o desenho mostra o caminho do simulador
        qualquer que seja, os testes daqui usam uma polia só e nada aqui diz
        onde fica a junção; o item ``The unwound sweep (PHY-54)'' do ADR-0004
        é editado aqui e no CLEAN-25, e o conflito de texto no rebase o stage 3
        resolve, porque o ADR está nos dois Primary.
  \item Um ticket, não dois --- umas vinte linhas de produção, uma fatia
        vertical do simulador até a tela; o simulador expondo \texttt{path}
        sozinho não muda nada que o usuário veja.
  \item \texttt{Difficulty: normal}, \texttt{Review: agent} --- nada numérico
        novo; o ponto que um primeiro passe pode errar, a leitura velha no
        editor, tem critério e mutação próprios.
  \item Sem critério próprio para a direção do PHY-45 no desenho --- vem da
        mesma linha que o critério 1 testa; um cenário a mais mediria o PHY-45
        de novo, não esta costura.
\end{enumerate}

\subsection*{PHY-55 (9)}
\begin{enumerate}
  \item O grip solto junta as duas peças numa só, de comprimento somado, e
        fica nela como polia ideal --- a regra do PHY-54 só no grip dá de
        +42,4\,kJ a +213,7\,kJ; juntar dá $\leq 0{,}0$\,J sem mudar a corda
        cujo grip segura.
  \item Soltar também pela previsão do fim do passo, num sentido só --- só na
        leitura real dá até +45,9\,kJ, e a previsão segue sem escrever
        \texttt{keep} nem \texttt{sweeps}.
  \item Histórico em todas as polias de uma corda com grip --- só no grip, a
        corda mista ganha até +5,3\,kJ; custa uma linha em \texttt{ropeFrame}.
  \item Cortar em dois, soltura aqui e reengate no PHY-57 --- a corda mista
        não pede código a mais, e o reengate tem critérios próprios que nenhum
        critério de energia separa.
  \item O critério 2 fica com os cenários da corda mista e bloqueia por
        CLEAN-25 --- 7 deles só passam com a regra de junções do CLEAN-25, e
        são os casos em que a corda sai das duas polias.
  \item O A 1/1 com \texttt{vx = 8} sai do critério 2 --- passa hoje
        ($-0{,}14$\,J) e não fica vermelho com nenhuma das mutações do
        critério 7.
  \item \texttt{Blocked by: CLEAN-25, CLEAN-26, PHY-56} --- o CLEAN-26 muda
        \texttt{gripShares} e o PHY-56 guarda o caminho em \texttt{ropeFrame},
        as duas funções que este ticket muda; o PHY-56 não diz nada sobre grip
        que conflite.
  \item O disco solto segue com o $\omega$ que tinha, e o critério de energia
        conta só os blocos --- o disco parte do repouso, então blocos acima de
        $E_0$ já é ganho do sistema.
  \item Limiar de 0,5\,J (o do PHY-54); \texttt{Difficulty: hard},
        \texttt{Review: agent} --- critérios que interagem e lógica numérica em
        funções que todo grip usa; nada de gosto.
\end{enumerate}

\subsection*{PHY-57 (5)}
\begin{enumerate}
  \item Cortar do PHY-55 --- com o grip solto para sempre a energia já fecha
        ($\leq +0{,}08$\,J), e só os critérios deste ticket separam o
        reengate.
  \item Prender de volta só com a peça juntada não frouxa --- uma corda frouxa
        não aperta o disco, e repartir a folga deixa a corda puxando até
        0,35\,m mais curta que $L$; esperar dá $\leq 0{,}046$\,m.
  \item O impacto do reengate sem código próprio --- dividir no comprimento
        atual deixa a diferença de velocidade para a previsão e a correção, que
        já tratam a corda que estica, de forma inelástica (energia
        $\leq +0{,}08$\,J).
  \item A marca no meio do arco, como em \texttt{buildWorld} --- manter a
        parte antiga dá os mesmos números, porque cada peça fica com o seu
        comprimento no caminho; não vira critério.
  \item Limiares 0,05\,N (medido $\leq 0{,}0054$\,N, mutante
        $\geq 10{,}69$\,N) e 0,1\,m (medido $\leq 0{,}046$\,m, mutante
        $\geq 0{,}28$\,m), janelas depois do reengate e antes de \texttt{a}
        chegar ao disco; \texttt{Difficulty: normal}, \texttt{Review: agent}
        --- mudança local de uma dúzia de linhas, com o protótipo medido.
\end{enumerate}
}

\section{Débito humano}

Nenhuma linha \texttt{Débito humano:} nos seis tickets, e nenhum
\texttt{Proxy escalated}; as 31 decisões tomadas pelo Bruno estão na seção
anterior. O PR \href{https://github.com/Laginho/PhysiSyst/pull/13}{\#13} tinha
as seis linhas como \texttt{Review: agent} --- Approve.

\paragraph{Depois do relatório (2026-10-02).} O Bruno mergeou o PR \#13 e
aprovou todas as recomendações. As 1--3 entraram no
\href{https://github.com/Laginho/Agent-Skills/pull/3}{Agent-Skills\#3}: o driver
grava \texttt{claude-opus-5-5} como \texttt{opus-5.5}; o scoreboard lê as
grafias antigas do mesmo jeito e conta numa coluna nova os achados que a
revisão consertou (rodada 0); e o foreman passa a mesclar o run log e a citar
as decisões do proxy no relatório. O log local foi sincronizado com a cópia
(178 linhas). Os três consertos da revisão viraram linhas de rodada 0,
rotuladas pelo GPT-6 Astra, e as seções 4 e 5 já os mostram. A 5 e a 6 viraram o CLEAN-28 (caminho da corda escolhido num
lugar só, mais o \texttt{drawScene} com 8 argumentos) e o CLEAN-29 (testes das
três regras do CLEAN-25), os dois à espera do stage 1. A 4 já estava feita, e a
7 não pede nada.

\end{document}
